\chapter{The Tower of Oracle Extensions}
\label{sec:fibration}
\label{ch:tower}
% ===================================================================

\section{Background: Turing's Oracle Hierarchy}

Turing~\cite{turing1939} described a transfinite tower of
computational systems above the Turing-computable.
At the base is the class of Turing-computable functions.
At the first stage above is the class of functions computable
relative to the \emph{halting oracle} $\emptyset'$: the oracle that
decides, for any Turing machine $M$ and input $x$, whether $M$
halts on $x$.
At the next stage is the class computable relative to
$\emptyset''$, the halting oracle for machines with access to
$\emptyset'$.
The tower continues through all countable ordinals and beyond.

Each stage is strictly more powerful than the stage below:
there are functions computable at stage $\alpha+1$ that are not
computable at stage $\alpha$, and this gap cannot be closed by
any finite amount of additional computation at stage $\alpha$.

\section{The Subcategory of Turing-Complete GSLTs}

\begin{definition}[Turing-Complete GSLT]
A GSLT $\GSLT$ is \emph{Turing-complete} if the class of functions
computable by terms of $\GSLT$ (via some standard encoding of
natural numbers) coincides with the class of Turing-computable
functions.
The full subcategory of Turing-complete GSLTs is written
$\mathbf{GSLT}_{\mathrm{TC}} \subseteq \mathbf{GSLT}$.
\end{definition}

\begin{example}
The $\lambda$-calculus, the $\pi$-calculus, and the Rho calculus
are all Turing-complete.
The category $\mathbf{GSLT}_{\mathrm{TC}}$ is therefore large and
contains the primary examples of interest.
\end{example}

\section{Oracle Extension of a GSLT}

\begin{definition}[Oracle Term]
Let $\GSLT \in \mathbf{GSLT}_{\mathrm{TC}}$.
An \emph{oracle term} for $\GSLT$ is a distinguished term
$\oracle \in \terms(\GSLT^{(1)})$ in an extension of $\GSLT$,
equipped with rewrite rules of the form:
\[
  \oracle \inter \lceil (M, x) \rceil \;\rewrite\;
  \begin{cases}
    \lceil \mathtt{halt} \rceil & \text{if } M \text{ halts on } x \\
    \lceil \mathtt{loop} \rceil & \text{otherwise}
  \end{cases}
\]
where $\lceil \cdot \rceil$ denotes the encoding of Turing machine
descriptions and inputs as terms of $\GSLT$, and $\inter$ is the
interaction operator of $\GSLT$.
The oracle answers halting queries in one interaction step.
\end{definition}

\begin{remark}[The Oracle as a Communicating Process]
In the Rho calculus, the oracle is naturally expressed as a
persistent receive on a distinguished channel $x_{\oracle}$:
\[
  \oracle \;=\; \mathtt{for}(q \leftarrow x_\oracle)\,
  \bigl( \text{compute halting of } {*}q \text{ and reply} \bigr)
\]
This makes the oracle a \emph{communicating process} rather than
a magic tape: it participates in the same interaction mechanism as
every other term.
This is essential if each stage of the tower is to be a GSLT rather
than something that escapes the GSLT framework at each oracle step.
\end{remark}

\begin{definition}[Oracle Extension]
\label{def:oracle-ext}
Let $\GSLT \in \mathbf{GSLT}_{\mathrm{TC}}$.
The \emph{first oracle extension} $\fiber{1}$ is the GSLT obtained
from $\GSLT$ by:
\begin{enumerate}[label=(\roman*)]
  \item extending the grammar with the oracle term $\oracle$;
  \item adding the oracle rewrite rules to $\rules$;
  \item extending the equations $\eqs$ to account for oracle
        interactions.
\end{enumerate}
The \emph{$\alpha$-th oracle extension} $\fiber{\alpha}$ is defined
by transfinite induction:
\begin{itemize}[itemsep=2pt]
  \item $\fiber{0} = \GSLT$;
  \item $\fiber{\alpha+1}$ is the oracle extension of $\fiber{\alpha}$,
        with oracle deciding halting in $\fiber{\alpha}$;
  \item $\fiber{\lambda} = \varinjlim_{\alpha < \lambda} \fiber{\alpha}$
        for limit ordinals $\lambda$ (colimit in $\mathbf{GSLT}$).
\end{itemize}
\end{definition}

\section{The Tower}

\begin{definition}[The oracle tower]
\label{def:oracle-tower}
The \emph{oracle tower} over $\mathbf{GSLT}_{\mathrm{TC}}$ is the
assignment
\[
  \Phi : \mathbf{GSLT}_{\mathrm{TC}} \times \Ord \;\to\; \mathbf{GSLT}
\]
sending $(\GSLT, \alpha)$ to $\fiber{\alpha}$, together with:
\begin{itemize}[itemsep=2pt]
  \item projection morphisms $\proj{\alpha} : \fiber{\alpha+1} \to \fiber{\alpha}$
        in $\mathbf{GSLT}$ for each ordinal $\alpha$;
  \item inclusion morphisms $\iota_\alpha : \fiber{\alpha} \to \fiber{\alpha+1}$
        embedding each stage into the next.
\end{itemize}
The tower over $\GSLT$ is
$\fiber{0} \hookrightarrow \fiber{1} \hookrightarrow \fiber{2}
\hookrightarrow \cdots$
\end{definition}

\begin{remark}[Why this is not called a fibration]
\label{rmk:not-a-fibration}
Earlier drafts called $\Phi$ the \emph{hypercomputational fibration},
and it is not one.
A Grothendieck fibration is a functor $p : \mathcal{E} \to \mathcal{B}$
equipped with cartesian lifts.
What Definition~\ref{def:oracle-tower} writes down is a functor
\emph{out of} a product --- an indexed family, from which the
Grothendieck construction would build a fibration, and which is
therefore the raw material for one rather than an instance of one.

Worse, the variance is overdetermined.
$\Ord$ is a poset, so a functor out of it can go one way, and both
$\proj{\alpha}$ and $\iota_\alpha$ are given.
What is being described is a section--retraction pair, and the
intended reading --- inclusion adds structure, projection forgets it
--- is that $\iota_\alpha \dashv \proj{\alpha}$.
If that adjunction holds, and the triangle identities have not been
checked here, the whole tower is a tower of adjunctions, which is a
better object than the one on the page and costs nothing that is not
already available.

The object this act is really about is the fibration
$p : \Hyp \to \catGSLT$ of open problem~(ii) in
Section~\ref{ndo:sec:open}, where the base is \emph{all} of
$\catGSLT$ rather than its Turing-complete part, and where the
question of whether angelic and demonic resolution are the Lawvere
adjoints $\exists \dashv (-)^{\ast} \dashv \forall$ to reindexing can
even be asked.
That object has not been exhibited.
This one has, and it is a ladder, so it is called one.
\end{remark}

\begin{proposition}[Projections are Morphisms]
Each projection $\proj{\alpha} : \fiber{\alpha+1} \to \fiber{\alpha}$
is a morphism in $\mathbf{GSLT}$: it preserves bisimulation.
If $P \bisim_{\alpha+1} Q$ in $\fiber{\alpha+1}$, then
$\proj{\alpha}(P) \bisim_\alpha \proj{\alpha}(Q)$ in $\fiber{\alpha}$.
\end{proposition}

\begin{proof}[Proof sketch]
Any bisimulation in $\fiber{\alpha+1}$ that uses only oracle-free
transitions is already a bisimulation in $\fiber{\alpha}$.
The projection discards oracle interactions; the remaining
bisimulation relation is preserved.
The details require checking that the oracle rewrite rules are
projected consistently with the extension of $\eqs$.
\end{proof}

\begin{remark}[Projections are Lossy]
The projection $\proj{\alpha}$ is in general \emph{not} a
bisimulation equivalence: there exist $P, Q$ with
$P \bisim_{\alpha+1} Q$ but $P \not\bisim_\alpha Q$.
Going down the tower collapses distinctions.
Going up reveals them.
\end{remark}



\section{Two things this construction gets wrong}
\label{sec:tower-wrong}

It is better to say this here than to let a reader discover it three
chapters on, so: the construction just given is a scaffold, and it has
two defects that the chapters after it exist to repair.

\paragraph{The index should not be an ordinal.}
Ordinals are linearly ordered, and the thing being indexed is not.
Two agents can each be able to resolve what the other cannot --- one
affords a choice principle the other lacks, and the other affords a
different one --- and an ordinal index has no room for that.
The correct index is a position in the Weihrauch lattice
\cite{weihrauch}, which has joins and meets and incomparable elements,
and the tower is the image of that lattice under a map that flattens it
into a chain.
A chain-shaped shadow is a legitimate way to introduce an intuition and
an illegitimate way to keep one.
Chapter~\ref{ch:apr} replaces it.

\paragraph{The oracle should not be bolted on.}
Definition~\ref{def:oracle-ext} adjoins an oracle term to a theory that
did not have one, with rewrite rules supplied by hand.
That is a construction and it is unobjectionable as a construction, but
it gets the situation backwards.
In any calculus where interaction is primitive and confluence fails,
something already decides which of several admissible pairings occurs,
and that something is not derivable from the terms.
It is the scheduler, and it is present at every cut.
The excess capacity we are looking for does not have to be added
to the theory; it has to be \emph{found} in it, and then priced.
Chapters~\ref{ch:apr} and~\ref{ch:agy} do the finding, and
Chapter~\ref{ch:chs} supplies the typing discipline that says what a
scheduler is allowed to be.

\paragraph{What survives.}
Both repairs leave the tower's actual content intact.
Projections are lossy, going up reveals distinctions, and each stage
carries a strictly richer physics than the one below --- those are
claims about oracle extensions and they do not depend on the index being
well-ordered or on the oracle being adjoined rather than discovered.
What does not survive is the picture of consciousness as a ladder with
rungs and a reader wondering which one they are standing on.

% ===================================================================
